\chapter{Differential Calculus and Taylor Approximations}

\section{Derivatives and Differentiability}

\subsection{The Derivative: Definition, Geometry, and Notation}

Let $f : D \to \mathbb{R}$ be a function defined in a neighborhood of a point $a \in \mathbb{R}$. The central object of differential calculus is the derivative of $f$ at $a$, which measures the instantaneous rate of change of $f$ near that point.

\begin{definition}
The derivative of a function $f$ at a point $a$, denoted by $f'(a)$, is defined as
\[
\begin{aligned}
f'(a) &= \lim_{h \to 0} \frac{f(a+h) - f(a)}{h} \\
&= \lim_{x \to a} \frac{f(x) - f(a)}{x - a},
\end{aligned}
\]
provided that this limit exists and is finite.
\end{definition}

For any given $h \in \mathbb{R}$, it is convenient to introduce the notation
\[
R_a(h) = \frac{f(a+h) - f(a)}{h}
\]
for the ratio appearing in the definition above, so that
\[
f'(a) = \lim_{h \to 0} R_a(h),
\]
provided the limit exists and is finite. The derivative admits several complementary interpretations, the first of which is geometric. Given a curve $C$ with equation $y = f(x)$, one seeks the tangent line to $C$ at the point $P = (a, f(a))$. To construct this tangent line, one considers a nearby point $Q = (x, f(x))$ with $x \neq a$, and computes the slope of the secant line through $P$ and $Q$,
\[
m_Q = \frac{f(x) - f(a)}{x - a}.
\]

Letting $Q$ approach $P$ along the curve, that is, letting $x$ approach $a$, one defines the tangent line $t$ to be the limiting position of the secant line, whenever such a limiting position exists.

\begin{definition}
The tangent line to the curve $C$ at the point $P$ is the line through $P$ with slope
\[
m = \lim_{x \to a} \frac{f(x) - f(a)}{x - a},
\]
provided that this limit exists and is finite.
\end{definition}

This limit is precisely the derivative $f'(a)$, so that the derivative may be identified with the slope of the tangent line to the graph of $f$ at the point $(a, f(a))$.

A second interpretation of the derivative arises from physics. Suppose an object moves along a straight line according to an equation of motion $s = f(t)$, where $s$ denotes the displacement of the object from the origin at time $t$. In order to determine the instantaneous velocity of the object at a given time $t = a$, one first considers the time interval from $t = a$ to $t = a + h$, over which the change in position is $f(a+h) - f(a)$; the average velocity over this interval is
\[
v_{\text{mean}} = \frac{\text{displacement}}{\text{time}} = \frac{f(a+h) - f(a)}{h}.
\]

Computing the average velocity over shorter and shorter time intervals, one defines the instantaneous velocity at time $t = a$ as the limit of these average velocities,
\[
v(a) = \lim_{h \to 0} \frac{f(a+h) - f(a)}{h},
\]
which again is recognized as the derivative $f'(a)$. A third and more general interpretation concerns rates of change. Suppose $y$ is a quantity depending on another quantity $x$, so that $y$ is a function of $x$ and one writes $y = f(x)$. If $x$ changes from $x_1$ to $x_2$, the change in $x$, also called the increment of $x$, is $\Delta x = x_2 - x_1$, and the corresponding change in $y$ is
\[
\Delta y = f(x_2) - f(x_1).
\]
The difference quotient, or incremental ratio,
\[
\frac{\Delta y}{\Delta x} = \frac{f(x_2) - f(x_1)}{x_2 - x_1}
\]
is called the average rate of change of $y$ with respect to $x$ over the interval $[x_1, x_2]$. Considering the average rate of change over smaller and smaller intervals, by letting $\Delta x$ approach zero, one defines the instantaneous rate of change of $y$ with respect to $x$ at $x_1$ as
\[
\lim_{\Delta x \to 0} \frac{\Delta y}{\Delta x} = \lim_{\Delta x \to 0} \frac{f(x_1 + \Delta x) - f(x_1)}{\Delta x},
\]
which is once again recognized as the derivative $f'(x_1)$. This general notion of instantaneous rate of change has broad applicability outside physics; for instance, biomedical scientists have studied the physiological changes in the body resulting from alcohol consumption by means of the law
\[
C(t) = 0.0225 \, t \, e^{-0.0467 t}
\]
milligrams per milliliter, and one may ask how quickly the blood alcohol concentration is increasing after ten minutes, a question that is naturally answered by evaluating the derivative $C'(10)$.

Several notations for the derivative coexist in the literature, namely $f'(x)$, $y'$, $df/dx$, $dy/dx$, $(d/dx) f(x)$, and $Df(x)$. Among these, the symbol $dy/dx$, known as Leibniz notation and dating back to 1672, deserves a specific comment: it should not be regarded as an actual ratio of two quantities $dy$ and $dx$, but simply as a synonym for $f'(x)$. Nonetheless, this notation is extremely useful and suggestive, since it directly invokes the definition of the derivative as a limit of a ratio of increments,
\[
\frac{dy}{dx} = \lim_{\Delta x \to 0} \frac{\Delta y}{\Delta x}.
\]

Figure~\ref{fig:atlas-secant-tangent} shows how secant slopes approach the derivative at a fixed point.

\begin{figure}[!htbp]
\centering
\begin{tikzpicture}[>=Stealth]
\begin{axis}[width=11cm,height=5.8cm,axis lines=middle,ticklabel style={font=\small},label style={font=\small},legend style={font=\scriptsize,draw=black!20,fill=white},xlabel={$x$},ylabel={$y$},xmin=0,xmax=2.4,ymin=-0.3,ymax=5,xtick={1,1.5,2},ytick={1,2,3,4}]
\addplot[figblue,very thick,domain=0:2.2,samples=100] {x^2};
\addplot[figred,dashed,thick,domain=0.6:2.15] {3*x-2};
\addplot[figred!55,dashdotted,thick,domain=0.6:2.15] {2.5*x-1.5};
\addplot[figgreen,thick,domain=0.4:2.1] {2*x-1};
\addplot[only marks,mark=*,figblue] coordinates {(1,1)(1.5,2.25)(2,4)};
\node[below right] at (axis cs:1,1) {$(1,1)$};
\end{axis}
\end{tikzpicture}
\caption{For $f(x)=x^2$ at $x=1$, the secants through the second points $x=2$ and $x=1.5$ have slopes $3$ and $2.5$. They approach the tangent slope $2$ as the second point tends to the first. The tangent is drawn in green.}
\label{fig:atlas-secant-tangent}
\end{figure}

\subsection{Differentiability and Continuity}

\begin{definition}
A function $f$ is differentiable at $a$ if $f'(a)$ exists. We say that $f$ is differentiable on an interval $I$ if it is differentiable at every point of $I$.
\end{definition}

Differentiability is a strictly stronger requirement than continuity, as recorded in the following theorem.

\begin{theorem}
If $f$ is differentiable at $a$, then $f$ is continuous at $a$.
\end{theorem}

\begin{remark}
The converse of this theorem is not true: the function $f(x) = |x|$ is continuous but not differentiable at $x = 0$.
\end{remark}

\begin{proof}
We wish to prove that
\[
\lim_{x \to a} f(x) = f(a),
\]
which is equivalent to showing that
\[
\lim_{x \to a} [f(x) - f(a)] = 0.
\]
Writing
\[
f(x) - f(a) = \frac{f(x) - f(a)}{x - a} \cdot (x-a)
\]
and passing to the limit as $x \to a$, using the differentiability of $f$ at $a$, we obtain
\[
\begin{aligned}
\lim_{x \to a} [f(x) - f(a)] &= \lim_{x \to a} \frac{f(x) - f(a)}{x-a} \cdot \lim_{x \to a} (x-a) \\
&= f'(a) \cdot 0 \\
&= 0,
\end{aligned}
\]
which proves the claim.
\end{proof}

There are essentially four distinct ways in which a function can fail to be differentiable at a point $x = a$. The most familiar is the presence of an angle in the graph, as exemplified by $f(x) = |x|$ at the origin: here $f$ is continuous at $a$, and the right-derivative and left-derivative, respectively defined by
\[
f'_+(a) := \lim_{h \to 0^+} \frac{f(a+h) - f(a)}{h}, \qquad f'_-(a) := \lim_{h \to 0^-} \frac{f(a+h) - f(a)}{h},
\]
both exist and are finite, but differ from one another. A second way is the presence of a discontinuity of $f$ itself at $a$, which automatically precludes differentiability. A third way is the presence of a vertical tangent or a cusp, in which case $f$ is continuous at $a$ but the right-derivative and left-derivative are both infinite.

\subsection{Derivative Functions and Higher Derivatives}

Given any number $x$ for which the limit
\[
f'(x) = \lim_{h \to 0} [f(x+h) - f(x)]/h
\]
exists and is finite, one may assign to $x$ the corresponding number $f'(x)$, obtaining a rule $x \mapsto f'(x)$. In this way, $f'$ may itself be regarded as a new function, called the derivative of $f$. The domain of $f'$ is the set
\[
\{x \in \operatorname{Dom}(f) : f'(x) \text{ is finite}\},
\]
and this set may be strictly smaller than the domain of $f$ itself. The graph of $f'$ can be derived from the graph of $f$ by exploiting the geometric fact that $f'(x)$ represents the slope of the tangent line to the graph of $f$ at the point $(x, f(x))$.


If $y = f(x)$ is a differentiable function, then its derivative $y' = f'(x)$ is itself a function of $x$, and it may happen that $f'$ possesses a derivative of its own, denoted by $(f')' = f''$. This new function is called the second derivative of $f$, since it is obtained by differentiating the derivative of $f$. Using Leibniz notation, the second derivative of $f$ is written as
\[
\begin{aligned}
y'' &= f''(x) \\
&= \frac{d}{dx}\left( \frac{dy}{dx} \right) \\
&= \frac{d^2 y}{dx^2}.
\end{aligned}
\]

A particularly natural physical interpretation of the second derivative is as a rate of change of a rate of change, the most familiar instance of which is acceleration. If $s = s(t)$ is the position function of an object moving along a straight line, its first derivative represents the velocity of the object as a function of time,
\[
\begin{aligned}
v(t) &= s'(t) \\
&= ds/dt;
\end{aligned}
\]
the instantaneous rate of change of velocity with respect to time is then called the acceleration of the object,
\[
\begin{aligned}
a(t) &= v'(t) \\
&= \frac{d^2 s}{dt^2}.
\end{aligned}
\]

The process of differentiation can be iterated further. The third derivative is defined as the derivative of the second derivative, $(f'')' = f'''$, and is written in Leibniz notation as
\[
\begin{aligned}
y''' &= f'''(x) \\
&= \frac{d}{dx}\left( \frac{d^2 y}{dx^2} \right) \\
&= \frac{d^3 y}{dx^3}.
\end{aligned}
\]

More generally, the $n$-th derivative of $f$ is denoted by $f^{(n)}$ and is obtained from $f$ by differentiating $n$ times in succession,
\[
\begin{aligned}
y^{(n)} &= f^{(n)}(x) \\
&= \frac{d^n y}{dx^n}.
\end{aligned}
\]

As will be seen in later developments of the theory, second and higher derivatives play a fundamental role in representing functions as infinite sums of power functions.

\section{Differentiation Rules}

\subsection{General Rules for Combining Derivatives}

When new functions are formed from previously known functions by means of addition, subtraction, multiplication, division, or composition, their derivatives can be computed directly in terms of the derivatives of the original functions. The theory of differentiation is governed by a small number of such rules, namely the rule for the sum, difference, and multiplication by a constant, known collectively as linearity, together with the product rule, the quotient rule, the chain rule for composition, and the rule for the derivative of the inverse function.

\begin{theorem}[Linearity]
Let $f$ and $g$ be differentiable at the point $x$, and let $c \in \mathbb{R}$. Then the functions $f + g$ and $c \cdot f$ are differentiable at $x$, and
\[
(f+g)'(x) = f'(x) + g'(x), \qquad (c \cdot f)'(x) = c \cdot f'(x).
\]

\end{theorem}

These rules follow immediately by applying the algebraic limit laws to the corresponding incremental ratios.

\begin{theorem}[Product Rule]
Let $f$ and $g$ be differentiable at the point $x$. Then the product function $f \cdot g$ is differentiable at $x$, and
\[
(f \cdot g)'(x) = f'(x) \cdot g(x) + f(x) \cdot g'(x).
\]

\end{theorem}

\begin{theorem}[Quotient Rule]
Let $f$ and $g$ be differentiable at the point $x$. If $g(x) \neq 0$, then the quotient $f/g$ is differentiable at $x$, and
\[
\left( \frac{f}{g} \right)'(x) = \frac{f'(x) \cdot g(x) - f(x) \cdot g'(x)}{[g(x)]^2}.
\]

\end{theorem}

\begin{theorem}[Chain Rule]
Let $g$ be differentiable at the point $x$, and let $f$ be differentiable at the point $y = g(x)$. Then the composition $f \circ g$ is differentiable at $x$, and
\[
(f \circ g)'(x) = f'(g(x)) \cdot g'(x).
\]

\end{theorem}

\begin{theorem}[Derivative of the Inverse]
Let $f : (a,b) \to \mathbb{R}$ be continuous and injective. If $f$ is differentiable at the point $c$ and $f'(c) \neq 0$, then the inverse function $f^{-1}$ is differentiable at the point $d = f(c)$, and
\[
(f^{-1})'(d) = \frac{1}{f'(c)}.
\]

\end{theorem}

We conclude this subsection with the proof of the product rule, since it illustrates the standard technique of adding and subtracting an auxiliary term in order to isolate the incremental ratios of the individual factors.

\begin{proof}[Proof of the Product Rule]
We must compute the limit of the incremental ratio of $f \cdot g$,
\[
R_h = \frac{f(x+h) g(x+h) - f(x) g(x)}{h}.
\]
Adding and subtracting the term $f(x+h) g(x)$, we obtain
\[
f(x+h) g(x+h) - f(x) g(x) = f(x+h) [g(x+h) - g(x)] + g(x) [f(x+h) - f(x)],
\]
so that
\[
R_h = f(x+h) \, \frac{g(x+h) - g(x)}{h} + g(x) \, \frac{f(x+h) - f(x)}{h}.
\]

In light of the differentiability of $f$ and $g$ at $x$, we may now pass to the limit as $h \to 0$. Observing that $f(x+h) \to f(x)$ by the continuity of $f$ at $x$, which is itself a consequence of the differentiability of $f$ at $x$, and that
\[
\frac{g(x+h) - g(x)}{h} \to g'(x), \qquad \frac{f(x+h) - f(x)}{h} \to f'(x),
\]
we conclude that
\[
\lim_{h \to 0} R_h = f(x) g'(x) + g(x) f'(x),
\]
which is precisely the statement of the product rule.
\end{proof}

\subsection{Derivatives of the Elementary Functions}

Combining the general differentiation rules with the definition of derivative, one can compute explicitly the derivative of each of the elementary functions cataloged earlier. Here the parameter $\alpha$ denotes an arbitrary real number, and every derivative formula stated below is understood to hold at every point $x$ where the right-hand side is well defined. We begin with the exponential function. Its derivative is given by
\[
\frac{d}{dx} e^x = e^x.
\]
Indeed, writing the incremental ratio and factoring out $e^x$, we have
\[
\lim_{h \to 0} \frac{e^{x+h} - e^x}{h} = e^x \lim_{h \to 0} \frac{e^h - 1}{h},
\]
and the result follows from the special limit
\[
\lim_{h \to 0} (e^h - 1)/h = 1
\]
established in the theory of infinitesimal sequences. The derivative of the natural logarithm is given by
\[
\frac{d}{dx} \ln x = \frac{1}{x}.
\]
Indeed, we compute
\[
\begin{aligned}
\lim_{h \to 0} \frac{\ln(x+h) - \ln x}{h} &= \lim_{h \to 0} \frac{\ln\!\left( \frac{x+h}{x} \right)}{h} \\
&= \frac{1}{x} \lim_{h \to 0} \frac{\ln\!\left(1 + \frac{h}{x}\right)}{\frac{h}{x}},
\end{aligned}
\]
and the conclusion again follows from the corresponding special limit for the logarithm. The derivatives of the sine and cosine functions are given by
\[
\frac{d}{dx} \sin x = \cos x, \qquad \frac{d}{dx} \cos x = -\sin x.
\]
To establish the first identity, we must compute
\[
\lim_{h \to 0} \frac{\sin(x+h) - \sin x}{h}.
\]
Using the addition formula
\[
\sin(x+h) = \sin x \cos h + \cos x \sin h,
\]
we obtain
\[
\lim_{h \to 0} \frac{\sin x \cos h + \cos x \sin h - \sin x}{h} = \sin x \cdot \lim_{h \to 0} \frac{\cos h - 1}{h} + \cos x \cdot \lim_{h \to 0} \frac{\sin h}{h}.
\]
Exploiting the special limits
\[
\begin{gathered}
\lim_{h \to 0} (\sin h)/h = 1 \\
\lim_{h \to 0} \frac{\cos h - 1}{h} = \lim_{h \to 0} (-h) \, \frac{1 - \cos h}{h^2} = 0,
\end{gathered}
\]
we conclude that the derivative of $\sin x$ equals $\cos x$; the derivative of $\cos x$ is obtained by a completely analogous computation. The derivative of a real power function is given by
\[
\frac{d}{dx} x^\alpha = \alpha \, x^{\alpha - 1}, \qquad \alpha \in \mathbb{R}.
\]
Indeed, one rewrites the increment $(x+h)^\alpha - x^\alpha$ as $x^\alpha [ ( 1 + h/x )^\alpha - 1 ]$, so that
\[
\frac{(x+h)^\alpha - x^\alpha}{h} = x^{\alpha - 1} \, \frac{\left(1 + \frac{h}{x}\right)^\alpha - 1}{\frac{h}{x}},
\]
and hence
\[
\begin{aligned}
\lim_{h \to 0} \frac{(x+h)^\alpha - x^\alpha}{h} &= x^{\alpha - 1} \lim_{h \to 0} \frac{\left(1 + \frac{h}{x}\right)^\alpha - 1}{\frac{h}{x}} \\
&= x^{\alpha - 1} \cdot \alpha,
\end{aligned}
\]
where the last limit is again computed by means of the corresponding special limit for infinitesimal sequences.

Using the product rule, one easily obtains the derivative of $\tan x$ from the derivatives of $\sin x$ and $\cos x$; using instead the rule for the derivative of the inverse function, one proves the derivatives of the inverse trigonometric functions,
\[
D \arcsin x = \frac{1}{\sqrt{1 - x^2}}, \qquad D \arccos x = -\frac{1}{\sqrt{1 - x^2}}, \qquad D \arctan x = \frac{1}{1 + x^2}.
\]

\begin{proof}
Let $f(x) = \sin x$. By definition, $y = \arcsin x$ if and only if $x = \sin y$. Applying the rule for the derivative of the inverse function, we obtain
\[
\begin{aligned}
D f^{-1}(x) &= \frac{1}{f'(y)} \\
&= \frac{1}{\cos y} \\
&= \frac{1}{\sqrt{1 - \sin^2 y}} \\
&= \frac{1}{\sqrt{1 - x^2}},
\end{aligned}
\]
where the positive square root is selected since $\cos y \geq 0$ on the range of $\arcsin$. The formula for $\arccos x$ is obtained analogously. For $f(x) = \tan x$, one instead has $y = \arctan x$ if and only if $x = \tan y$, and hence
\[
\begin{aligned}
D f^{-1}(x) &= \frac{1}{f'(y)} \\
&= \frac{1}{\frac{1}{\cos^2 y}} \\
&= \cos^2 y \\
&= \frac{\cos^2 y}{\cos^2 y + \sin^2 y} \\
&= \frac{1}{1 + \tan^2 y} \\
&= \frac{1}{1 + x^2},
\end{aligned}
\]
which completes the proof.
\end{proof}

\section{The Main Theorems of Differential Calculus}

Having established the basic computational rules for derivatives, we now turn to three fundamental theoretical results, namely Fermat's Theorem, Rolle's Theorem, and Lagrange's Theorem, together with several applications to optimization, the behavior of graphs, and the computation of limits.

\subsection{Fermat's Theorem, Critical Points, and Optimization}

We begin by making precise the notions of local extremum that underlie the theory. Let $f$ be defined in a neighborhood of a point $c \in \mathbb{R}$.

\begin{definition}
The point $c$ is a local minimum point for $f$ if there exists a neighborhood $U(c)$ such that $f(x) \geq f(c)$ for every $x \in U(c)$, and $c$ is a local maximum point for $f$ if there exists a neighborhood $U(c)$ such that $f(x) \leq f(c)$ for every $x \in U(c)$. In either case, the corresponding value $f(c)$ is called a local minimum, respectively local maximum, value for $f$.
\end{definition}

The following theorem provides a necessary condition, expressed in terms of the derivative, for the existence of a local extremum.

\begin{theorem}[Fermat's Theorem]
If $f$ has a local maximum or a local minimum at $c$, and $f$ is differentiable at $c$, then $f'(c) = 0$.
\end{theorem}

\begin{definition}
A point $c$ in the domain of $f$ such that $f'(c) = 0$ is called a critical point of $f$.
\end{definition}

Fermat's Theorem may thus be rephrased as follows: if $f$ has a local maximum or minimum at $c$, then either $f$ is not differentiable at $c$, or $c$ is a critical point of $f$.

\begin{proof}
Suppose, for the sake of definiteness, that $f$ has a local maximum at $c$; the case of a local minimum is treated analogously. Then $f(c) \geq f(x)$ for every $x$ sufficiently close to $c$, which implies that for $h$ sufficiently close to $0$,
\begin{equation*}
f(c) \geq f(c+h) \implies f(c+h) - f(c) \leq 0. \tag{$*$}
\end{equation*}
If $h > 0$ is sufficiently small, dividing this inequality by $h$ preserves its direction, giving
\[
\frac{f(c+h) - f(c)}{h} \leq 0.
\]

Taking the right-hand limit as $h \to 0^+$ and invoking the Sign Preserving Property for limits, we obtain
\[
\lim_{h \to 0^+} \frac{f(c+h) - f(c)}{h} \leq 0 \implies f'_+(c) \leq 0.
\]

Analogously, if $h < 0$ is sufficiently small, dividing inequality $(*)$ by $h$ reverses its direction, giving
\[
\frac{f(c+h) - f(c)}{h} \geq 0,
\]
so that taking the left-hand limit yields
\[
\lim_{h \to 0^-} \frac{f(c+h) - f(c)}{h} \geq 0 \implies f'_-(c) \geq 0.
\]
Since $f$ is differentiable at $c$, we have
\[
\begin{aligned}
f'_+(c) &= f'_-(c) \\
&= f'(c),
\end{aligned}
\]
and the only value compatible with the two inequalities $f'(c) \leq 0$ and $f'(c) \geq 0$ is $f'(c) = 0$, as claimed.
\end{proof}


Fermat's Theorem provides the theoretical foundation for a systematic method of finding the absolute extreme values of a function on a closed bounded interval. Let $f$ be a continuous function on a closed bounded interval $[a,b]$; the problem is to find the absolute maximum and minimum values of $f$ on $[a,b]$, that is, two points $c, d \in [a,b]$ such that
\[
\begin{aligned}
f(c) &\leq f(x) \\
&\leq f(d)
\end{aligned}
\]
for every $x \in [a,b]$. Weierstrass' Theorem guarantees that such points exist, but it does not, by itself, indicate how to locate them; this is precisely the role played by Fermat's Theorem.

The resulting procedure, known as the Closed Interval Method, consists of three steps. The first step is to find the critical points of $f$ inside the open interval $(a,b)$, including any points where $f$ fails to be differentiable, and to evaluate $f$ at each of these points. The second step is to evaluate $f$ at the two endpoints of the interval, namely at $a$ and at $b$. The third step is to compare all the values obtained in the first two steps: the largest of these values is the absolute maximum value of $f$ on $[a,b]$, and the smallest is the absolute minimum value.

As an illustration of this method, one may consider the function $g(x) = x^4 - 4x^3$ on the interval $[0,5]$, and ask what its maximum and minimum values are on this interval, as well as what happens on the open interval $(0,5)$ and on the whole real line $\mathbb{R}$. A further classical application concerns an open box constructed from a sheet of metal measuring twenty centimetres by twelve centimetres, from each corner of which a square of side length $x$ is removed before the sides are folded upward; the problem is to determine the value of $x$ that maximizes the resulting volume of the box.

\subsection{Rolle's Theorem and Lagrange's Theorem}

A second fundamental theorem, known as the Mean Value Theorem, generalizes Fermat's Theorem by relating the derivative of a function at an interior point to the average rate of change of the function over an entire interval.

\begin{theorem}[Lagrange's Theorem]
Let $f$ be a function satisfying the following two hypotheses: $f$ is continuous on the closed interval $[a,b]$, and $f$ is differentiable on the open interval $(a,b)$. Then there exists a number $c \in (a,b)$ such that
\[
f'(c) = \frac{f(b) - f(a)}{b - a}.
\]

\end{theorem}

A particularly important special case of Lagrange's Theorem, in which the values of $f$ at the two endpoints coincide, is traditionally stated as a separate result.

\begin{theorem}[Rolle's Theorem]
Let $f$ be a function satisfying the following three hypotheses: $f$ is continuous on the closed interval $[a,b]$, $f$ is differentiable on the open interval $(a,b)$, and $f(a) = f(b)$. Then there exists a number $c \in (a,b)$ such that $f'(c) = 0$.
\end{theorem}

\begin{proof}
The proof distinguishes three cases according to the behavior of $f$ on $(a,b)$. In the first case, $f$ is constant on $[a,b]$; then $f' \equiv 0$ throughout, so the number $c$ may be taken to be any point of $(a,b)$. In the second case, $f$ is not constant and there exists some $x \in (a,b)$ such that
\[
f(x) > f(a) = f(b);
\]
by Weierstrass' Theorem, $f$ attains a maximum value somewhere on $[a,b]$, and since this maximum value exceeds $f(a) = f(b)$, it must be attained at a point $c$ inside the open interval $(a,b)$, so that $c$ is a local maximum of $f$. Since $f$ is differentiable at $c$, Fermat's Theorem yields $f'(c) = 0$. In the third case, $f$ is not constant and there exists some $x \in (a,b)$ such that
\[
f(x) < f(a) = f(b);
\]
by Weierstrass' Theorem, $f$ attains a minimum value somewhere on $[a,b]$, and since this minimum value is smaller than $f(a) = f(b)$, it must be attained at a point $c$ inside $(a,b)$. Again, Fermat's Theorem yields $f'(c) = 0$, which completes the proof in all three cases.
\end{proof}

\begin{proof}[Proof of Lagrange's Theorem]
We define the auxiliary function
\[
F(x) = f(x) - \frac{f(b) - f(a)}{b - a} (x - a).
\]

Observe that $F$ is continuous on $[a,b]$ and differentiable on $(a,b)$, being a linear combination of $f$ and of an affine function of $x$. Moreover, a direct computation shows that
\[
\begin{aligned}
F(a) &= f(a) \\
&= F(b),
\end{aligned}
\]
so that $F$ satisfies the hypotheses of Rolle's Theorem on $[a,b]$. Applying Rolle's Theorem to $F$, we find a point $c \in (a,b)$ such that $F'(c) = 0$. For every $x \in (a,b)$, the derivative of $F$ is
\[
F'(x) = f'(x) - \frac{f(b) - f(a)}{b-a}.
\]
Setting $x = c$ yields
\[
0 = F'(c) = f'(c) - \frac{f(b) - f(a)}{b-a} \implies f'(c) = \frac{f(b) - f(a)}{b-a},
\]
which establishes the theorem.
\end{proof}

Figure~\ref{fig:atlas-mean-value} illustrates the parallel chord and tangent in the Mean Value Theorem.

\begin{figure}[!htbp]
\centering
\begin{tikzpicture}[>=Stealth]
\begin{axis}[width=11cm,height=5.8cm,axis lines=middle,ticklabel style={font=\small},label style={font=\small},legend style={font=\scriptsize,draw=black!20,fill=white},xlabel={$x$},ylabel={$f(x)$},xmin=-0.15,xmax=2.3,ymin=-0.4,ymax=4.7,xtick={0,2},ytick={1,2,3,4}]
\addplot[figblue,very thick,domain=0:2.15,samples=100] {x^2};
\addplot[figred,thick] coordinates {(0,0)(2,4)};
\addplot[figgreen,thick,domain=0.5:2] {2*x-1};
\addplot[only marks,mark=*,figblue] coordinates {(0,0)(1,1)(2,4)};
\draw[dotted] (axis cs:1,0) -- (axis cs:1,1);
\node[below] at (axis cs:1,0) {$c$};
\end{axis}
\end{tikzpicture}
\caption{For $f(x)=x^2$ on $[0,2]$, the endpoint chord has slope $2$. The tangent at $c=1$ has the same slope, as required by the Mean Value Theorem. The red chord records average change; the green tangent records instantaneous change.}
\label{fig:atlas-mean-value}
\end{figure}

\section{L'H\^opital's Rule and Its Applications}

\subsection{The Rule and Its Proof}

Having established the main theorems of differential calculus, we now turn to a natural question: how can derivatives be used to compute limits, particularly limits presenting an indeterminate form? The following rule, based ultimately on Lagrange's Theorem and a generalization thereof due to Cauchy, provides a powerful tool for this purpose.

\begin{theorem}[L'H\^opital's Rule]
Let $c \in \mathbb{R}$, and let $f, g$ be differentiable in a punctured neighborhood $\dot{U}(c)$ of $c$. Assume that the following three conditions hold: first,
\[
\begin{aligned}
\lim_{x \to c} f(x) &= \lim_{x \to c} g(x) \\
&= 0;
\end{aligned}
\]
second, $g'(x) \neq 0$ throughout $\dot{U}(c)$; and third,
\[
\lim_{x \to c} f'(x)/g'(x) = \ell
\]
for some $\ell \in \mathbb{R}$. Then
\[
\lim_{x \to c} \frac{f(x)}{g(x)} = \ell.
\]

\end{theorem}

\begin{remark}
The rule remains valid for indeterminate forms of the type $\infty / \infty$, for one-sided limits, and for limits taken as $x \to \pm \infty$.
\end{remark}

The proof of L'H\^opital's Rule rests on a generalization of Lagrange's Theorem due to Cauchy, which we now state.

\begin{theorem}[Cauchy's Theorem]
Let $f, g$ be continuous on $[a,b]$ and differentiable on $(a,b)$, with $g' \neq 0$ on $(a,b)$. Then there exists $c \in (a,b)$ such that
\[
\frac{f(b) - f(a)}{g(b) - g(a)} = \frac{f'(c)}{g'(c)}.
\]

\end{theorem}

\begin{proof}[Proof of L'H\^opital's Rule]
We present a short proof under the additional assumption that $f$ and $g$ are continuous at $c$, so that, in particular,
\[
\begin{aligned}
f(c) &= g(c) \\
&= 0.
\end{aligned}
\]

Fix any $x > c$ sufficiently close to $c$. Since $g' \neq 0$ on $(c,x)$, Cauchy's Theorem may be applied to the interval $[c,x]$, yielding a point $\xi_x \in (c,x)$ such that
\[
\begin{aligned}
\frac{f(x)}{g(x)} &= \frac{f(x) - f(c)}{g(x) - g(c)} \\
&= \frac{f'(\xi_x)}{g'(\xi_x)}.
\end{aligned}
\]

As $x \to c$, one has $\xi_x \to c$ as well, since $\xi_x$ lies between $c$ and $x$; the third hypothesis of the theorem then guarantees that the right-hand side tends to $\ell$, and the conclusion follows.
\end{proof}

\subsection{Repeated Application and Indeterminate Forms}

As an illustration of the rule, consider the limit
\[
\lim_{x \to 0} (\sin x - x)/x^3.
\]

Taking the quotient of the derivatives of numerator and denominator, one obtains $(\cos x - 1)/(3x^2)$; applying L'H\^opital's Rule once more to this new quotient yields $-\sin x / (6x)$, and one further application yields $-\cos x / 6$. In summary, in light of L'H\^opital's Rule applied three times in succession,
\[
\begin{aligned}
\lim_{x \to 0} \frac{\sin x - x}{x^3}
&= \lim_{x \to 0} \frac{\cos x - 1}{3x^2} \\
&= \lim_{x \to 0} \frac{-\sin x}{6x}
= \lim_{x \to 0} \frac{-\cos x}{6}
= -\frac{1}{6}.
\end{aligned}
\]

It is worth noting that the computation may in fact be stopped one step earlier, at the quotient $-\sin x / (6x)$, since the value $-1/6$ then follows directly from the special limit
\[
\lim_{x \to 0} (\sin x)/x = 1
\]
established previously, without a further application of the rule. L'H\^opital's Rule is particularly useful for establishing the hierarchy of infinities as $x \to +\infty$, namely the relation $(\ln x)^\beta \prec x^\alpha \prec e^x$ for $\alpha > 0$; consider, for instance, the comparison $\ln x \prec x^\alpha$. Indeed, for $\alpha > 0$, applying L'H\^opital's Rule gives
\[
\lim_{x \to +\infty} \frac{\ln x}{x^\alpha} \xrightarrow{\text{l'H\^opital}} \lim_{x \to +\infty} \frac{\frac{1}{x}}{\alpha x^{\alpha - 1}} = \lim_{x \to +\infty} \frac{1}{\alpha x^\alpha} = 0.
\]

The rule is likewise useful for resolving indeterminate forms of type $0 \cdot \infty$; for instance, the special limit
\[
\lim_{x \to 0^+} x^\alpha (\ln x)^\beta = 0,
\]
in the particular case $\beta = 1$, can be established by rewriting $x^\alpha \ln x$ as the quotient $\ln x / x^{-\alpha}$, which presents the indeterminate form $-\infty / \infty$, and applying L'H\^opital's Rule:
\[
\begin{aligned}
\lim_{x \to 0^+} x^\alpha \ln x
&= \lim_{x \to 0^+} \frac{\ln x}{x^{-\alpha}} \\
&\xrightarrow{\text{l'H\^opital}} \lim_{x \to 0^+} \frac{\frac{1}{x}}{-\alpha x^{-\alpha - 1}} \\
&= \lim_{x \to 0^+} \left( -\frac{x^\alpha}{\alpha} \right) = 0.
\end{aligned}
\]

Despite its evident usefulness, L'H\^opital's Rule must be applied with a degree of caution, and it is important to be aware of its limitations. In particular, the rule cannot be used to prove the special limits from which it itself indirectly draws, since some of those limits are logically prior to the derivatives of the elementary functions involved. Moreover, the rule cannot be applied to limits such as
\[
\begin{gathered}
\lim_{x \to 1^+} x/(x-1) \\
\lim_{x \to +\infty} (x + \sin x)/x,
\end{gathered}
\]
since in the first case the hypothesis of a $0/0$ or $\infty/\infty$ indeterminate form is not satisfied, while in the second case the derivative of the denominator does not lead to a resolvable limit in the required sense. Finally, there exist situations in which the rule is simply inconclusive, such as the limit
\[
\lim_{x \to \infty} (x^2+1)^{1/2}/x,
\]
where repeated differentiation of numerator and denominator fails to terminate in a determinate expression, so that alternative methods must be used instead.

\section{Monotonicity, Concavity, and Graph Analysis}

\subsection{Monotonicity}

The sign of the first derivative of a function is directly related to its monotone behavior, as expressed by the following test, which is itself a consequence of Lagrange's Theorem.

\begin{theorem}[Increasing/Decreasing Test]
Let $f$ be differentiable on an interval $I$. If $f'(x) = 0$ for every $x \in I$, then $f$ is constant on $I$. If $f' > 0$ on $I$, then $f$ is increasing on $I$. If $f' < 0$ on $I$, then $f$ is decreasing on $I$.
\end{theorem}

\begin{proof}
Let $x_1 < x_2$ be any two numbers in $I$. Applying Lagrange's Theorem to the interval $[x_1, x_2]$, we find a number $c$ between $x_1$ and $x_2$ such that
\[
f(x_2) - f(x_1) = f'(c)(x_2 - x_1).
\]

If $f' > 0$ on $I$, this identity immediately shows that $f(x_2) - f(x_1) > 0$, proving that $f$ is increasing on $I$; the remaining two cases are established by an entirely analogous argument.
\end{proof}

In view of this test, it follows that a local maximum or minimum can occur at any point $c$ where the first derivative changes sign, and this observation is made precise in the following criterion.

\begin{theorem}[First Derivative Test]
Assume that $f$ is continuous at a point $c$ and differentiable in a punctured neighborhood $\dot{U}(c)$. If $f'$ changes sign from positive to negative at $c$, then $f$ has a local maximum at $c$. If $f'$ changes sign from negative to positive at $c$, then $f$ has a local minimum at $c$. If $f'$ does not change sign at $c$, then $f$ has neither a local maximum nor a local minimum at $c$.
\end{theorem}

\begin{remark}
It should not be assumed that all local minima look essentially alike upon close inspection of the graph: there exist local minima at which the function fails to be decreasing on the left and increasing on the right of the point in question, so that the First Derivative Test, while sufficient, does not capture every possible local extremum.
\end{remark}

A further useful consequence of Lagrange's Theorem provides a sufficient condition for differentiability at a point, based on the limiting behavior of the derivative nearby.

\begin{theorem}
Assume $f$ is continuous at a point $c$ and differentiable in a punctured neighborhood $\dot{U}(c)$. If
\[
\lim_{x \to c} f'(x) = L \in \mathbb{R},
\]
then $f$ is differentiable at $c$ and $f'(c) = L$.
\end{theorem}

\begin{proof}
By Lagrange's Theorem applied to the interval $[c, c+h]$, for $h$ sufficiently small, there exists a point $x_h \in (c, c+h)$ such that
\[
\frac{f(c+h) - f(c)}{h} = f'(x_h).
\]
Since $x_h \to c$ as $h \to 0$, it follows that
\[
\begin{aligned}
f'_+(c) &= \lim_{h \to 0^+} \frac{f(c+h) - f(c)}{h} \\
&= \lim_{h \to 0^+} f'(x_h) \\
&= L,
\end{aligned}
\]
and an analogous argument for $h \to 0^-$ shows that $f'_-(c) = L$ as well, so that $f'(c) = L$, as claimed.
\end{proof}

\begin{remark}
The value $L = \infty$ is also admissible in this theorem: in that case one cannot conclude that $f$ is differentiable at $c$, but the value $f'(c) = \infty$ is naturally interpreted as indicating the presence of a vertical tangent to the graph at the point $(c, f(c))$, as illustrated by the functions
\[
\begin{gathered}
f(x) = x^{1/3} \\
f(x) = \sqrt[3]{x^2(6-x)}.
\end{gathered}
\]

It is also worth noting that, as a consequence of the Darboux Property, the derivative of a function, regarded as a function in its own right, can never exhibit a jump discontinuity or a removable discontinuity.
\end{remark}

\subsection{Concavity and Inflection Points}

Beyond monotonicity, the shape of a graph is further characterized by its concavity, which describes whether the graph lies above or below its tangent lines.

\begin{definition}
Given a function $f$ that is differentiable on an interval $I$, we say that $f$ is convex, or concave upward, on $I$ if, for every $c \in I$, the graph of $f$ lies above the tangent line to the graph at the point $(c, f(c))$. Equivalently, convexity is characterized by the following inequality, required for every $x, c \in I$:
\[
f \text{ convex on } I \iff f(x) \geq f(c) + f'(c)(x-c).
\]

We say that $f$ is concave, or concave downward, on $I$ whenever the opposite inequality holds throughout $I$.
\end{definition}

\begin{definition}
A point $P = (c, f(c))$ on a curve $y = f(x)$ is called an inflection point if $f$ is continuous at $c$ and the curve changes from concave upward to concave downward, or from concave downward to concave upward, at $c$.
\end{definition}

It is worth noting that at an inflection point where $f$ is differentiable, the curve crosses its own tangent line, in contrast with the behavior at a generic point of the graph. The relationship between concavity and the monotonicity of the derivative is made precise by the following theorem.

\begin{theorem}
Let $f$ be differentiable on an interval $I$. Then the following two statements are equivalent: first, $f$ is convex on $I$; second, $f'$ is increasing on $I$.
\end{theorem}

\begin{proof}
We first show that the second statement implies the first. Take $c \in I$ and $x > c$. By Lagrange's Theorem applied to the interval $[c,x]$, there exists $y \in (c,x)$ such that
\[
\frac{f(x) - f(c)}{x - c} = f'(y).
\]
Since $f'$ is increasing, we have $f'(y) \geq f'(c)$, and hence
\[
\begin{aligned}
f(x) - f(c) &= f'(y)(x-c) \\
&\geq f'(c)(x-c),
\end{aligned}
\]
so that
\[
f(x) \geq f'(c)(x-c) + f(c),
\]
which is precisely the definition of convexity at $c$, from the right; the case $x < c$ is treated analogously, establishing convexity on all of $I$.

We now show that the first statement implies the second. Let $x < y$ be two points of $I$. By the definition of convexity, the value $f(y)$ is greater than the value of the tangent line to the graph at $(x, f(x))$ evaluated at $y$, that is,
\[
f(y) \geq f(x) + f'(x)(y-x), \qquad \text{hence} \qquad f'(x) \leq \frac{f(y) - f(x)}{y - x}.
\]

Similarly, the value $f(x)$ is greater than the value of the tangent line to the graph at $(y, f(y))$ evaluated at $x$, that is,
\[
f(x) \geq f(y) + f'(y)(x-y), \qquad \text{hence} \qquad f'(y) \geq \frac{f(y) - f(x)}{y - x}.
\]

Combining the two inequalities, we immediately deduce that $f'(x) \leq f'(y)$, which shows that $f'$ is increasing on $I$ and completes the proof of the equivalence.
\end{proof}

As a direct consequence of this equivalence, one obtains a practical criterion for concavity expressed in terms of the second derivative.

\begin{theorem}[Concavity Test]
Let $f$ be twice differentiable on an interval $I$. If $f'' > 0$ on $I$, then $f$ is convex on $I$. If $f'' < 0$ on $I$, then $f$ is concave, or concave downward, on $I$.
\end{theorem}

In view of this test, it follows that a point of inflection can occur at any point where the second derivative changes sign; more precisely, inflection points are to be sought among those points $c$ at which either $f''(c) = 0$ or $f$ fails to be twice differentiable. These principles, together with the Increasing/Decreasing Test, the First Derivative Test, and the Closed Interval Method discussed earlier, provide a complete set of tools for sketching the graph of a given function, as may be illustrated by sketching the graphs of
\[
\begin{gathered}
g(x) = x^4 - 4x^3 \\
f(x) = \sqrt[3]{x^2(6-x)}.
\end{gathered}
\]

\section{Taylor Approximation and Asymptotic Expansions}

\subsection{Local Polynomial Approximation}

A fundamental geometric fact about differentiable functions is that, upon zooming in toward a point on their graph, the graph increasingly resembles its tangent line at that point. This observation suggests that the tangent line provides a natural first-order approximation to a differentiable function near a given point, and it is this idea that we now develop rigorously. Let $f$ be differentiable at a point $a$, and consider the tangent line approximation
\[
L(x) = f(a) + f'(a)(x-a).
\]

This linear function is, in a precise sense, the best linear approximation to $y = f(x)$ near $x = a$: among all first-degree polynomials of the form $P(x) = c_0 + c_1 x$, the function $L(x)$ is the one that fits the graph of $f$ most closely near the point $(a, f(a))$. This distinguished status is characterized algebraically by the fact that $L(x)$ is the unique first-degree polynomial satisfying the two conditions $P(a) = f(a)$ and $P'(a) = f'(a)$, that is, the unique linear polynomial that agrees with $f$ both in value and in first derivative at the point $a$.

The precise sense in which $L(x)$ approximates $f(x)$ is clarified by the following consequence of the very definition of differentiability.

\begin{theorem}
Let $f$ be differentiable at $a$. Then
\[
\lim_{x \to a} \frac{f(x) - L(x)}{x - a} = 0.
\]

\end{theorem}

This result may be interpreted as saying that the error
\[
E(x) = f(x) - L(x)
\]
approaches zero faster than $(x-a)$ itself as $x \to a$; adopting the standard Landau notation, we write $E(x) = o(x-a)$ as $x \to a$. With this notation, the conclusion of the theorem may be restated as
\[
f(x) = f(a) + f'(a)(x-a) + o(x-a), \qquad x \to a,
\]
or, equivalently, setting $x = a + h$,
\[
f(a+h) = f(a) + f'(a) h + o(h), \qquad h \to 0.
\]

Having established that the tangent line provides a first-order approximation to $f$ near $a$, it is natural to ask whether a better approximation can be obtained by allowing a higher-degree polynomial. Specifically, one may attempt to approximate the curve $y = f(x)$ near $(a, f(a))$ not by a straight line but by a parabola, that is, by a second-degree, or quadratic, approximation of the form
\[
P(x) = c_0 + c_1 x + c_2 x^2.
\]

A concrete numerical experiment with the function $y = \cos x$ centered at $a = 0$ illustrates vividly how such a quadratic approximation improves upon the linear one.

Inspired by this example, we search for a polynomial $P(x)$ satisfying the three conditions $P(a) = f(a)$, $P'(a) = f'(a)$, and
\[
P''(a) = f''(a).
\]

It turns out that there is a unique quadratic function satisfying these three conditions simultaneously, namely
\[
T_2(x) = f(a) + f'(a)(x-a) + \frac{1}{2} f''(a) (x-a)^2.
\]

This function is called the second-degree Taylor polynomial of $f$ centered at $a$, and it is customarily denoted by $T_2(x)$. A graphical comparison, for instance between the function $f(x) = \cos x$ and its second-degree Taylor polynomial centered at $a = 0$, makes it apparent that the quadratic approximation is considerably more accurate than the linear one over a wider range of values of $x$; this observation is not merely qualitative but reflects a precise general fact, stated in the following theorem.

\begin{theorem}[Taylor--Peano, Degree Two]
If $f$ is twice differentiable at $a$, then
\[
\lim_{x \to a} \frac{f(x) - T_2(x)}{(x-a)^2} = 0.
\]

\end{theorem}

As in the linear case, this result may be interpreted in terms of the error term: the error
\[
E_2(x) = f(x) - T_2(x)
\]
approaches zero faster than $(x-a)^2$ as $x \to a$, that is,
\[
E_2(x) = o((x-a)^2)
\]
as $x \to a$. The conclusion of the theorem may accordingly be rewritten as
\[
f(x) = f(a) + f'(a)(x-a) + \frac{1}{2} f''(a)(x-a)^2 + o\big((x-a)^2\big), \qquad x \to a,
\]
or, setting $x = a+h$,
\[
f(a+h) = f(a) + f'(a) h + \frac{1}{2} f''(a) h^2 + o(h^2), \qquad h \to 0.
\]

\begin{proof}
We prove the result for $a = 0$; the general case follows by an entirely analogous argument after a translation of the variable. Setting
\[
g(x) = f(x) - T_2(x),
\]
the thesis becomes
\[
\lim_{x \to 0} g(x)/x^2 = 0.
\]
Observe that, by the very construction of $T_2$, we have
\[
\begin{aligned}
g(0) &= g'(0) \\
&= g''(0) \\
&= 0.
\end{aligned}
\]

As a first step, we apply Lagrange's Theorem to $g$ on the interval between $0$ and $x$: there exists $\tau \in (0,1)$ such that
\begin{equation*}
\frac{g(x) - g(0)}{x - 0} = g'(\tau x) \implies g(x) = g'(\tau x) \, x. \tag{1}
\end{equation*}
As a second step, we exploit the linear approximation of the function $g'$ at the point $0$, namely
\[
g'(y) = g'(0) + g''(0) y + o(y), \qquad y \to 0.
\]
Setting $y = \tau x$ and recalling that
\[
\begin{aligned}
g'(0) &= g''(0) \\
&= 0,
\end{aligned}
\]
we find
\begin{equation*}
g'(\tau x) = o(\tau x). \tag{2}
\end{equation*}
Combining $(1)$ and $(2)$, we obtain $g(x) = o(\tau x) \, x$, and hence
\[
\lim_{x \to 0} \frac{g(x)}{x^2} = \lim_{x \to 0} \frac{o(\tau x) \, x}{x^2} = \lim_{x \to 0} \frac{o(\tau x)}{\tau x} \cdot \tau = 0,
\]
which completes the proof.
\end{proof}

Figure~\ref{fig:ch5-taylor} compares a function with its linear and quadratic approximations near the expansion point.

\begin{figure}[!htbp]
\centering
\begin{tikzpicture}
  \begin{axis}[
    width=12.5cm,height=7.2cm,
    axis lines=middle,
    xlabel={$x$},ylabel={$y$},
    xmin=-2.6,xmax=2.6,ymin=-2.35,ymax=1.45,
    xtick={-2,0,2},xticklabels={$-2$,$a=0$,$2$},
    ytick={-2,-1,1},
    ticklabel style={font=\small},
    label style={font=\small},
    legend style={
      at={(0.5,1.03)},anchor=south,
      legend columns=3,
      font=\scriptsize,draw=black!30,fill=white,
      legend cell align=left,column sep=0.8em
    }
    ]
    \addplot[figblue,very thick,domain=-2.5:2.5,samples=180] {cos(deg(x))};
    \addlegendentry{$f(x)=\cos x$}
    \addplot[figred,thick,dashed,domain=-2.5:2.5] {1};
    \addlegendentry{$T_1(x)=1$}
    \addplot[figgreen,thick,dash dot,domain=-2.5:2.5,samples=180] {1-x^2/2};
    \addlegendentry{$T_2(x)=1-x^2/2$}
    \node[circle,fill=black,inner sep=1.25pt] at (axis cs:0,1) {};
  \end{axis}
\end{tikzpicture}
\caption[Linear and quadratic Taylor approximation]{Comparison between $f(x)=\cos x$, its
  first-degree Taylor polynomial $T_1(x)=1$, and its second-degree Taylor polynomial
  $T_2(x)=1-x^2/2$, all centered at $a=0$. The quadratic polynomial follows the curvature of $f$ and therefore
  remains accurate on a wider neighborhood of the expansion point.}
\label{fig:ch5-taylor}
\end{figure}

\subsection{Taylor's Theorem and Remainder Formulas}

The construction of the second-degree Taylor polynomial generalizes naturally to any degree $n$, provided that the function under consideration is differentiable a sufficient number of times.

\begin{definition}
Let $f$ be differentiable $n$ times at $a$. The polynomial
\[
T_n(x) := f(a) + f'(a)(x-a) + \frac{f''(a)}{2}(x-a)^2 + \cdots + \frac{f^{(n)}(a)}{n!}(x-a)^n
\]
is called the $n$-th degree Taylor polynomial of $f$ centered at $a$. In the special case $a = 0$, it is also called the $n$-th degree Maclaurin polynomial of $f$.
\end{definition}

By construction, the polynomial $T_n$ satisfies
\[
T_n^{(k)}(a) = f^{(k)}(a)
\]
for every $k = 1, \dots, n$, and it is in fact the unique polynomial of degree $n$ possessing this property. As the degree $n$ increases, the corresponding Taylor polynomial provides an increasingly accurate approximation to $f$ near $a$, a phenomenon that becomes visually apparent when several Taylor polynomials of increasing degree are plotted together against the original function.

This improvement in accuracy with increasing degree is made precise by the following generalization of the Taylor--Peano theorem stated above for the case $n=2$.

\begin{theorem}[Taylor--Peano Formula]
Let $f$ be differentiable $n$ times at $a$. Then
\[
\lim_{x \to a} \frac{f(x) - T_n(x)}{(x-a)^n} = 0.
\]

\end{theorem}

As before, this result may be phrased in terms of the error term: the error
\[
E_n(x) = f(x) - T_n(x)
\]
approaches zero faster than $(x-a)^n$ as $x \to a$, that is,
\[
E_n(x) = o((x-a)^n)
\]
as $x \to a$. The conclusion of the theorem is thus written as
\[
f(x) = f(a) + f'(a)(x-a) + \cdots + \frac{1}{n!} f^{(n)}(a)(x-a)^n + o\big((x-a)^n\big), \qquad x \to a,
\]
or, setting $x = a + h$,
\[
f(a+h) = f(a) + f'(a) h + \cdots + \frac{1}{n!} f^{(n)}(a) h^n + o(h^n), \qquad h \to 0.
\]

An instructive feature of the Taylor--Peano formula is that the remainder term $o(x^n)$ appearing at one order can itself be expanded further, revealing additional terms of the expansion at the next order; in this sense, one may zoom into the approximation as much as desired. For instance, as $x \to 0$ one has the first-order expansion $\sin x = x + o(x)$; zooming further into the remainder reveals that
\[
o(x) = -(1/6) x^3 + o(x^3),
\]
so that
\[
\sin x = x - \frac{1}{6} x^3 + o(x^3).
\]
Zooming in once again, the next term shows that
\[
o(x^3) = (1/5!) x^5 + o(x^5),
\]
yielding the still more refined expansion
\[
\sin x = x - \frac{1}{6} x^3 + \frac{1}{120} x^5 + o(x^5).
\]

This technique of successively refined Maclaurin expansions provides a powerful and systematic method for computing limits presenting an indeterminate form. As an illustration, consider the limit
\[
\lim_{x \to 0} (\sin x - x)/x^3.
\]
The Maclaurin expansion of $\sin x$ to third order gives
\[
\sin x = x - (1/6) x^3 + o(x^3),
\]
and hence
\[
\begin{aligned}
\lim_{x \to 0} \frac{\sin x - x}{x^3} &= \lim_{x \to 0} \frac{-\frac{1}{6} x^3 + o(x^3)}{x^3} \\
&= -\frac{1}{6} + \lim_{x \to 0} \frac{o(x^3)}{x^3}.
\end{aligned}
\]

Recalling the very definition of the symbol "small $o$", the second limit on the right-hand side vanishes by definition, and we conclude that
\[
\lim_{x \to 0} \frac{\sin x - x}{x^3} = -\frac{1}{6}.
\]

The Taylor--Peano formula, discussed above, describes the local behavior of a function near a point $a$, in the sense that it controls the error only asymptotically as $x \to a$, without providing any explicit bound valid for $x$ at a fixed, finite distance from $a$. A complementary result, known as the Taylor--Lagrange formula, instead provides an explicit expression for the remainder term, valid throughout an entire interval, and is therefore of a global rather than purely local character.

\begin{theorem}[Taylor--Lagrange Formula]
Let $f$ be differentiable $n+1$ times on an interval $[a,x]$. Then there exists a point $z \in (a,x)$ such that
\[
\begin{aligned}
f(x) &= T_n(x) + E_n(x), \\
E_n(x) &= \frac{f^{(n+1)}(z)}{(n+1)!} (x-a)^{n+1}.
\end{aligned}
\]

The expression for $E_n(x)$ given above is called the Lagrange form of the remainder term, in contrast with the Peano form
\[
E_n(x) = o((x-a)^n)
\]
furnished by the earlier theorem.
\end{theorem}

The global character of the Taylor--Lagrange formula becomes particularly valuable whenever an upper bound is available for the derivatives of $f$ throughout the interval in question, since in that case the Lagrange remainder can be estimated explicitly, yielding quantitative error bounds rather than merely qualitative asymptotic information. This feature is exploited decisively in the following application.

\subsection{Taylor Expansions of Elementary Functions}

It follows from the general theory developed above that, for every $n \in \mathbb{N}$,
\[
e^x = \sum_{k=0}^{n} \frac{x^k}{k!} + o(x^n), \qquad x \to 0,
\]
which shows that the Maclaurin polynomials of the exponential function reconstruct $e^x$ near $a = 0$ with arbitrary precision, in the local sense described by the Taylor--Peano formula. A far more striking fact, however, is that the Maclaurin polynomials of $e^x$ reconstruct the function not merely near $a=0$, but throughout its entire domain, in the sense made precise by the following theorem.

\begin{theorem}
For every $x \in \mathbb{R}$,
\[
e^x = \sum_{n=0}^{\infty} \frac{x^n}{n!}.
\]

\end{theorem}

Setting $x = 1$ in this identity yields the celebrated relation between Euler's number and the factorial series,
\[
e = \sum_{n=0}^{\infty} \frac{1}{n!}.
\]

\begin{proof}
Let $x > 0$ be fixed; the case $x \leq 0$ is handled by an entirely analogous argument. By the Taylor--Lagrange formula applied to $f(x) = e^x$ centered at $a = 0$, we have
\begin{equation*}
e^x - T_n(x) = \frac{f^{(n+1)}(z)}{(n+1)!} x^{n+1} \tag{$*$}
\end{equation*}
for some $z$ with $0 < z < x$, where $f^{(k)}(z) = e^z$ for every $k \geq 0$, and
\[
T_n(x) = \sum_{k=0}^{n} x^k / k!.
\]
Since $z < x$, we have $e^z < e^x$, and from $(*)$ we obtain the bound
\[
|e^x - T_n(x)| \leq e^x \, \frac{x^{n+1}}{(n+1)!}.
\]
Recalling that
\[
\lim_{n \to \infty} q^{n+1}/(n+1)! = 0
\]
for every fixed real number $q$, a fact established in the theory of sequences, we conclude that
\[
\lim_{n \to \infty} [e^x - T_n(x)] = 0,
\]
Equivalently, the limiting representation is
\[
\begin{aligned}
e^x &= \lim_{n \to \infty} T_n(x) \\
&= \sum_{n=0}^{\infty} \frac{x^n}{n!},
\end{aligned}
\]
where the last equality follows directly from the definition of series as the limit of its partial sums.
\end{proof}

The phenomenon just illustrated for the exponential function is an instance of a more general theory, that of Taylor series, or Maclaurin series in the case $a = 0$: for a sufficiently regular function $f$, the sequence of Maclaurin polynomials can reconstruct the entire function over its domain, not merely a local approximation near a single point. A useful intuitive image for this phenomenon is that the Taylor series of a function discovers, so to speak, the mathematical DNA encoded within the function at a single point, and allows the entire function to be rebuilt from this single piece of local information. This intuition is well illustrated by the sequence of Taylor polynomials of increasing degree for the function $y = \sin x$ centered at $a = 0$, which approximate the sine function over an ever wider range of values of $x$ as the degree increases.

The systematic use of Maclaurin expansions for the computation of limits, illustrated earlier for the limit
\[
\lim_{x \to 0} (\sin x - x)/x^3,
\]
rests on a small collection of algebraic manipulation rules governing the symbol $o(x^n)$. We now record the precise definition of this symbol, together with its basic algebraic properties.

\begin{definition}
Given $n \in \mathbb{N}$, the symbol $o(x^n)$, read "small $o$ of $x^n$", denotes a generic, unspecified quantity depending on $x$, vanishing as $x \to 0$, with the property that
\[
\lim_{x \to 0} \frac{o(x^n)}{x^n} = 0.
\]

\end{definition}

As $x \to 0$, the symbol $o(x^n)$ obeys the following algebraic rules, which together constitute what is customarily called the algebra of small $o$. First, $x^n = o(x^m)$ for any integer $m < n$, reflecting the fact that a higher power of $x$ is itself negligible compared to a lower power as $x \to 0$. Second, $c \cdot o(x^n) = o(x^n)$ for any constant $c$, so that multiplying a small-$o$ quantity by a fixed constant does not change its order. Third,
\[
o(x^n) \pm o(x^n) = o(x^n),
\]
so that the sum or difference of two quantities of the same order remains of that order. Fourth,
\[
x^n \cdot o(x^m) = o(x^{n+m}),
\]
so that multiplying a small-$o$ quantity by an explicit power increases its order accordingly. Fifth,
\[
o(x^n) \cdot o(x^m) = o(x^{n+m}),
\]
extending the previous rule to the product of two small-$o$ quantities. Sixth,
\[
(o(x^n))^m = o(x^{nm}),
\]
governing the behavior of small-$o$ quantities under exponentiation. Seventh,
\[
o(x^n + o(x^n)) = o(x^n),
\]
which shows that the symbol $o(x^n)$ is stable under composition with itself.

\subsection{Applications to Indeterminate Forms}

The algebra of small $o$ underlies a general method for computing limits of an indeterminate form $0/0$, of the type
\[
\lim_{x \to 0} f(x)/g(x).
\]

If $f$ and $g$ possess derivatives of every order at $x = 0$, the Maclaurin theorem allows one to write $f(x) = a x^n + o(x^n)$ for some integer $n$ and some nonzero constant $a$, meaning that $f$ is an infinitesimal of order $n$ with leading part $a x^n$, and analogously $g(x) = b x^m + o(x^m)$ for suitable $m$ and nonzero $b$. Substituting these expansions into the quotient, one obtains
\[
\begin{aligned}
\lim_{x \to 0} \frac{f(x)}{g(x)} &= \lim_{x \to 0} \frac{a x^n + o(x^n)}{b x^m + o(x^m)} \\
&= \lim_{x \to 0} \frac{x^n \left( a + \dfrac{o(x^n)}{x^n} \right)}{x^m \left( b + \dfrac{o(x^m)}{x^m} \right)},
\end{aligned}
\]
and it then suffices to apply the definition of small $o$, according to which the two fractions $o(x^n)/x^n$ and $o(x^m)/x^m$ both tend to zero as $x \to 0$, in order to resolve the limit explicitly in terms of $n$, $m$, $a$, and $b$.

\begin{remark}[Exercises on Taylor expansions and limits]
The techniques developed in this section admit a broad range of applications, several of which are proposed here as exercises for the reader. One may first show, for $x \to 0$ and any two integers $m < n$, that $x^n = o(x^m)$, and subsequently deduce the multiplicative rule
\[
x^n \cdot o(x^m) = o(x^{n+m})
\]
directly from the definition of small $o$. A second exercise consists of finding the Taylor polynomial of degree six for $f(x) = \sin x$ centered at $a = \pi/3$, thereby practicing the construction of Taylor polynomials at a center other than the origin. A third group of exercises, relying on the algebra of small $o$ together with the standard table of Maclaurin expansions, asks the reader to find the Maclaurin polynomial of degree four for $f(x) = \ln(\cos x)$; to determine the order of infinitesimal and the leading part, as $x \to 0$, of the function
\[
f(x) = (\sinh x)^2 - \sin(x^2);
\]
and to compute the limit
\[
\lim_{x \to 0} \frac{\tan x - \sin x}{x^2 (x+1) + \ln(1 - x^2)}
\]
by means of suitable Maclaurin expansions of numerator and denominator. A further set of exercises concerns the qualitative behavior of graphs near points where several derivatives vanish. One may first sketch a local graph of $f(x) = \ln(\cos x)$ close to $x = 0$, using the Maclaurin expansion obtained above. A second, more theoretical exercise concerns the classification of critical points by means of higher-order derivatives: suppose that
\[
\begin{gathered}
f'(a) = \cdots = f^{(n-1)}(a) = 0 \\
\begin{aligned}
f^{(n)}(a) &= k \\
&\neq 0
\end{aligned}
\end{gathered}
\]
for some integer $n \geq 2$; one is asked to prove that if $n$ is odd, then $a$ is an inflection point with horizontal tangent, while if $n$ is even, then $a$ is a local minimum when $k > 0$ and a local maximum when $k < 0$. A useful hint for this exercise, in the illustrative case $n = 2$, is that it suffices to know the values $f(a)$, $f'(a)$, and $f''(a)$: indeed, if $f''(a) > 0$, then as $x \to a$ one may write
\[
\begin{aligned}
f(x) &= f(a) + f'(a)(x-a) + \frac{f''(a)}{2}(x-a)^2 + o\big((x-a)^2\big) \\
&= f(a) + f'(a)(x-a)
+ (x-a)^2 \left[ \frac{f''(a)}{2} + \frac{o\big((x-a)^2\big)}{(x-a)^2} \right],
\end{aligned}
\]
and since the bracketed term is positive for $x$ sufficiently close to $a$, it follows that
\[
f(x) \geq f(a) + f'(a)(x-a)
\]
near $a$, that is, the graph of $f$ stays above its tangent line, confirming that $a$ is a local minimum when $n=2$ and $k>0$. The general case follows a similar pattern: for $h$ small, one has
\[
f(a+h) - f(a) = h^n \left[ \frac{k}{n!} + \frac{o(h^n)}{h^n} \right],
\]
and the sign of the bracketed expression, which approaches $k/n!$ as $h \to 0$, together with the sign of $h^n$, which depends on the parity of $n$, together determine the nature of the critical point.

A final group of exercises draws together the theory of Taylor series with the theory of series developed earlier. Using the table of Maclaurin expansions, one may find the sum of the series
\[
\sum_{n=0}^{\infty} (-1)^n x^{4n}/n!,
\]
for an arbitrary real number $x$, and of the numerical series
\[
\sum_{n=1}^{\infty} (-1)^{n-1} 3^n/(n \cdot 5^n),
\]
by recognizing each as a Maclaurin series evaluated at a suitable point. A second exercise asks the reader to use Maclaurin polynomials together with the Taylor--Lagrange Theorem to compute the numerical value of $e^{\pi/3}$ with a precision of $10^{-2}$; the standard approach is to approximate $e^{\pi/3}$ by the partial sum
\[
T_n(x) = \sum_{k=0}^{n} x^k/k!
\]
evaluated at $x = \pi/3$, choosing $n$ large enough that the Lagrange remainder satisfies $|E_n| < 10^{-2}$, using the bound
\[
0 < E_n(\pi/3) \leq \max_{z \in [0, \pi/3]} \frac{f^{(n+1)}(z)}{(n+1)!} \left( \frac{\pi}{3} \right)^{n+1} \leq e^{\pi/3} \, \frac{1}{(n+1)!} \left( \frac{\pi}{3} \right)^{n+1},
\]
and selecting the smallest integer $n$ for which the right-hand side falls below $10^{-2}$. A final, more theoretical exercise asks the reader to prove that Euler's number $e$ is irrational, a classical result that can be established using the series representation of $e$ together with a careful estimate of the remainder term of its Taylor--Lagrange expansion.
\end{remark}
